\documentclass[12pt, a4paper]{ctexart}
\usepackage{amsmath, amssymb, amsthm}
\usepackage{fancyhdr}
\usepackage{geometry}
\usepackage{enumitem}
\usepackage{graphicx}
\usepackage{tikz}
\usepackage{pgfplots}
\pgfplotsset{compat=1.18}

\geometry{left=2cm, right=2cm, top=2.5cm, bottom=2.5cm}

\setCJKmainfont{Hiragino Sans GB}
\setCJKsansfont{Heiti SC}

\pagestyle{fancy}
\fancyhf{}
\fancyhead[L]{高考数学满分突破·概率论公理化视角}
\fancyhead[R]{测试卷}
\fancyfoot[C]{\thepage}
\renewcommand{\headrulewidth}{0.4pt}

\setlength{\parindent}{0pt}
\setlength{\parskip}{6pt}

\newcommand{\score}[1]{\hfill\textbf{（#1分）}}

\begin{document}

\begin{center}
{\LARGE\bfseries 高考数学满分突破测试卷（六）}\\[8pt]
{\large 概率论公理化视角}\\[6pt]
考试时间：120分钟 \quad 满分：150分
\end{center}

\vspace{10pt}

{\large\bfseries 一、单项选择题：本题共8小题，每小题5分，共40分。每小题给出的四个选项中，只有一项是符合题目要求的。}

\vspace{8pt}

\begin{enumerate}

\item 设随机事件 $A$，$B$ 满足 $P(A)=0.4$，$P(B)=0.5$，$P(A\cup B)=0.7$，则 $P(A\mid B)=$ \score{5}
\begin{enumerate}[label=\Alph*.]
    \item $0.2$
    \item $0.4$
    \item $0.5$
    \item $0.6$
\end{enumerate}

\item 从公理化概率论的角度，设样本空间 $\Omega=\{1,2,3,4,5,6\}$，事件域 $\mathcal{F}$ 为 $\Omega$ 的全体子集构成的 $\sigma$-代数，概率测度 $P$ 满足 $P(\{k\})=\dfrac{k}{21}$，$k=1,2,\cdots,6$。则事件 $A=\{2,4,6\}$ 的概率 $P(A)=$ \score{5}
\begin{enumerate}[label=\Alph*.]
    \item $\dfrac{2}{7}$
    \item $\dfrac{4}{7}$
    \item $\dfrac{1}{2}$
    \item $\dfrac{3}{7}$
\end{enumerate}

\item 设 $X$ 服从参数为 $\lambda=2$ 的泊松分布，则 $P(X\geq 2)=$ \score{5}
\begin{enumerate}[label=\Alph*.]
    \item $1-3e^{-2}$
    \item $1-5e^{-2}$
    \item $3e^{-2}$
    \item $5e^{-2}$
\end{enumerate}

\item 甲、乙两人独立射击同一目标，甲命中率为 $0.6$，乙命中率为 $0.7$。已知目标被击中，则目标是被甲击中的概率为 \score{5}
\begin{enumerate}[label=\Alph*.]
    \item $\dfrac{3}{8}$
    \item $\dfrac{6}{13}$
    \item $\dfrac{7}{13}$
    \item $\dfrac{3}{7}$
\end{enumerate}

\item 从条件概率的链式法则出发，设 $P(A)=0.5$，$P(B\mid A)=0.6$，$P(C\mid A\cap B)=0.8$，则 $P(A\cap B\cap C)=$ \score{5}
\begin{enumerate}[label=\Alph*.]
    \item $0.24$
    \item $0.16$
    \item $0.30$
    \item $0.48$
\end{enumerate}

\item 设随机变量 $X$ 的分布列为 $P(X=k)=\dfrac{C}{k(k+1)}$，$k=1,2,3,\cdots$，其中 $C$ 为常数，则 $C=$ \score{5}
\begin{enumerate}[label=\Alph*.]
    \item $1$
    \item $2$
    \item $\dfrac{1}{2}$
    \item $\dfrac{1}{3}$
\end{enumerate}

\item 用概率的连续性（下连续性）解释：设 $A_1\subset A_2\subset A_3\subset\cdots$ 为递增事件列，$\displaystyle\bigcup_{n=1}^{\infty}A_n=A$，则 $\displaystyle\lim_{n\to\infty}P(A_n)=P(A)$。现设 $A_n=\left\{X>\dfrac{1}{n}\right\}$，$X$ 为 $(0,1)$ 上的均匀随机变量，则 $\displaystyle\lim_{n\to\infty}P(A_n)=$ \score{5}
\begin{enumerate}[label=\Alph*.]
    \item $0$
    \item $\dfrac{1}{2}$
    \item $1$
    \item 不存在
\end{enumerate}

\item 设 $X_1,X_2,\cdots,X_{100}$ 为独立同分布的随机变量，均服从 $B(1,0.3)$（参数为 $1,0.3$ 的伯努利分布），$\bar{X}=\dfrac{1}{100}\displaystyle\sum_{i=1}^{100}X_i$。由中心极限定理，$P(\bar{X}\leq 0.35)$ 的近似值为（已知 $\Phi(1.09)\approx 0.8621$） \score{5}
\begin{enumerate}[label=\Alph*.]
    \item $0.8621$
    \item $0.1379$
    \item $0.9000$
    \item $0.7500$
\end{enumerate}

\end{enumerate}

\vspace{10pt}

{\large\bfseries 二、多项选择题：本题共4小题，每小题5分，共20分。每小题有多个选项符合题目要求，全部选对得5分，部分选对得2分，有选错的得0分。}

\vspace{8pt}

\begin{enumerate}[resume]

\item 设 $A$，$B$ 为两个随机事件，且 $0<P(A)<1$，$0<P(B)<1$，则下列命题正确的是 \score{5}
\begin{enumerate}[label=\Alph*.]
    \item 若 $A$ 与 $B$ 互斥，则 $P(A\mid B)=0$
    \item 若 $A$ 与 $B$ 相互独立，则 $P(A\mid B)=P(A)$
    \item $P(A)=P(A\mid B)P(B)+P(A\mid \bar{B})P(\bar{B})$
    \item 若 $P(A\mid B)>P(A)$，则 $P(B\mid A)>P(B)$
\end{enumerate}

\item 关于大数定律，下列说法正确的是 \score{5}
\begin{enumerate}[label=\Alph*.]
    \item 切比雪夫大数定律要求随机变量序列两两不相关且方差有公共上界
    \item 辛钦大数定律要求随机变量独立同分布且数学期望存在
    \item 伯努利大数定律是辛钦大数定律的特殊情况
    \item 大数定律的结论是依概率收敛
\end{enumerate}

\item 设二维随机变量 $(X,Y)$ 的联合概率密度为 $f(x,y)=\begin{cases} 2, & 0<x<1,\, 0<y<x \\ 0, & \text{其他} \end{cases}$，则下列正确的是 \score{5}
\begin{enumerate}[label=\Alph*.]
    \item $f_X(x)=2x$，$0<x<1$
    \item $f_Y(y)=2(1-y)$，$0<y<1$
    \item $X$ 与 $Y$ 相互独立
    \item $P(X+Y<1)=\dfrac{1}{3}$
\end{enumerate}

\item 关于期望和方差的性质，下列正确的是 \score{5}
\begin{enumerate}[label=\Alph*.]
    \item $E(aX+b)=aE(X)+b$
    \item $D(aX+b)=a^2D(X)$
    \item 若 $X$，$Y$ 相互独立，则 $D(X-Y)=D(X)-D(Y)$
    \item $D(X)=E(X^2)-[E(X)]^2$
\end{enumerate}

\end{enumerate}

\vspace{10pt}

{\large\bfseries 三、填空题：本题共4小题，每小题5分，共20分。}

\vspace{8pt}

\begin{enumerate}[resume]

\item 设事件 $A$，$B$ 相互独立，$P(A)=0.3$，$P(A\cup B)=0.58$，则 $P(B)=$ \underline{\hspace{2cm}}。\score{5}

\item 设随机变量 $X\sim N(2, \sigma^2)$，若 $P(X<4)=0.84$，则 $P(0<X<2)=$ \underline{\hspace{2cm}}。\score{5}

\item 将一枚均匀硬币独立地抛掷 $n$ 次，设 $X$ 表示正面向上的次数。由De Moivre--Laplace中心极限定理，要使 $P\left(\left|\dfrac{X}{n}-\dfrac{1}{2}\right|<0.05\right)\geq 0.95$，则 $n$ 至少为 \underline{\hspace{2cm}}。（已知 $\Phi(1.96)\approx 0.975$）\score{5}

\item 设 $X$ 的概率密度为 $f(x)=\dfrac{1}{\sqrt{2\pi}}e^{-\frac{x^2}{2}}$（标准正态分布），则 $E(X^4)=$ \underline{\hspace{2cm}}。\score{5}

\end{enumerate}

\vspace{10pt}

{\large\bfseries 四、解答题：本题共6小题，共70分。解答应写出文字说明、证明过程或演算步骤。}

\vspace{8pt}

\begin{enumerate}[resume]

\item （10分）从公理化概率论的视角，设 $\Omega=\{a,b,c\}$，$\mathcal{F}=2^{\Omega}$，概率测度 $P$ 满足 $P(\{a\})=\alpha$，$P(\{b\})=\beta$，$P(\{c\})=\gamma$。

\begin{enumerate}[label=(\arabic*)]
    \item 写出概率公理要求 $\alpha,\beta,\gamma$ 满足的条件；
    \item 设 $A=\{a,b\}$，$B=\{b,c\}$，证明 $P(A\cap B)=P(A)+P(B)-P(A\cup B)$（容斥原理）；
    \item 若 $\alpha=\dfrac{1}{2}$，$\beta=\dfrac{1}{3}$，求 $\gamma$ 并计算 $P(A\mid B)$。
\end{enumerate}

\item （12分）甲、乙、丙三人参加某项测试，三人通过的概率分别为 $\dfrac{1}{2}$，$\dfrac{2}{3}$，$\dfrac{3}{4}$，且三人是否通过相互独立。

\begin{enumerate}[label=(\arabic*)]
    \item 求三人中恰有两人通过的概率；
    \item 设 $X$ 表示通过测试的人数，求 $X$ 的分布列和数学期望 $E(X)$；
    \item 已知三人中至少有两人通过，求甲通过的概率。
\end{enumerate}

\item （12分）设随机变量 $X$ 的概率密度函数为
$$f(x)=\begin{cases} ax+b, & 0<x<1 \\ 0, & \text{其他} \end{cases}$$
且 $E(X)=\dfrac{7}{12}$。

\begin{enumerate}[label=(\arabic*)]
    \item 求常数 $a$，$b$ 的值；
    \item 求 $D(X)$；
    \item 求 $P\left(X>\dfrac{1}{2}\right)$。
\end{enumerate}

\item （12分）设 $(X,Y)$ 为二维离散型随机变量，其联合分布列为：

\begin{center}
\begin{tabular}{c|ccc}
$X\backslash Y$ & $0$ & $1$ & $2$ \\ \hline
$0$ & $\dfrac{1}{6}$ & $\dfrac{1}{3}$ & $\dfrac{1}{12}$ \\
$1$ & $\dfrac{1}{6}$ & $0$ & $\dfrac{1}{4}$
\end{tabular}
\end{center}

\begin{enumerate}[label=(\arabic*)]
    \item 求边缘分布列 $P(X=0)$，$P(X=1)$；
    \item 判断 $X$ 与 $Y$ 是否独立，并说明理由；
    \item 求 $\text{Cov}(X,Y)$ 和相关系数 $\rho_{XY}$。
\end{enumerate}

\item （12分）用特征函数的方法研究二项分布的极限。

\begin{enumerate}[label=(\arabic*)]
    \item 设 $X\sim B(n,p)$，写出 $X$ 的特征函数 $\varphi_X(t)$；
    \item 当 $n\to\infty$，$p=\dfrac{\lambda}{n}$（$\lambda>0$ 为常数）时，证明 $\varphi_X(t)\to e^{\lambda(e^{it}-1)}$；
    \item 利用上述结论说明二项分布在此条件下收敛于参数为 $\lambda$ 的泊松分布。
\end{enumerate}

\item （12分）从信息论视角看概率分布。

\begin{enumerate}[label=(\arabic*)]
    \item 定义离散随机变量 $X$ 的熵为 $H(X)=-\displaystyle\sum_{k}p_k\ln p_k$，其中 $p_k=P(X=k)$。证明当 $X$ 服从两点分布 $B(1,p)$ 时，$H(X)=-p\ln p-(1-p)\ln(1-p)$；
    \item 证明：当 $X$ 取 $n$ 个值时，$H(X)\leq \ln n$，等号当且仅当 $X$ 服从均匀分布时成立（利用Jensen不等式或Lagrange乘数法）；
    \item 设 $X\sim B(1,0.5)$，$Y\sim B(1,0.9)$，分别计算 $H(X)$ 和 $H(Y)$，并比较哪个不确定性更大。
\end{enumerate}

\end{enumerate}

\end{document}
